\section{R ile çözümlü uygulama}
\label{sec:r-b03}

\textbf{Soru.} Çalışma süreleri 2, 3, 3, 4 ve 13 saat olan örnekte merkez,
yayılım ve aykırı değer sınırlarını bulun. Bölümde kullanılan doğrusal
enterpolasyonla uyum için yüzdelik hesabında \texttt{type = 7} seçin.

\begin{lstlisting}[style=prismR]
saat <- c(2, 3, 3, 4, 13)
ceyrek <- quantile(saat, c(0.25, 0.75), type = 7)
ceyrek_acikligi <- IQR(saat, type = 7)
c(ortalama = mean(saat), medyan = median(saat),
  varyans = var(saat), standart_sapma = sd(saat))
quantile(saat, c(0, 0.25, 0.5, 0.75, 1), type = 7)
sinir <- c(ceyrek[1] - 1.5 * ceyrek_acikligi,
           ceyrek[2] + 1.5 * ceyrek_acikligi)
sinir
saat[saat < sinir[1] | saat > sinir[2]]
hist(saat, breaks = c(0, 3, 6, 9, 12, 15), right = FALSE,
     main = "Calisma suresi", xlab = "Saat")
boxplot(saat, horizontal = TRUE, xlab = "Saat")
c(ortalama_13_haric = mean(saat[saat != 13]),
  medyan_13_haric = median(saat[saat != 13]))
\end{lstlisting}

\begin{outputguidebox}
\textbf{Çözüm ve kontrol değerleri.} Ortalama 5, medyan 3, örneklem
varyansı 20,5 ve örneklem standart sapması 4,5277'dir. Beş sayı özeti
$(2;3;3;4;13)$; çeyrekler açıklığı 1; alt ve üst sınırlar 1,5 ve 5,5'tir.
13 saat üst sınırı aşar. Bu değer çıkarıldığında ortalama 3'e iner;
medyan 3 olarak kalır.
\end{outputguidebox}

\textbf{Yorum.} \texttt{sd} ve \texttt{var} örneklem hesabında $n-1$
bölenini kullanır. \texttt{boxplot} kutu sınırları için Tukey menteşelerini
kullanır; bunlar bu örnekte seçilen çeyreklerle aynıdır, her örneklemde aynı
olmak zorunda değildir. Aykırı değer işareti otomatik silme gerekçesi değildir.

\textbf{Pekiştirme ve yanıt.} 13'ü çıkarma işlemi yalnız duyarlılık
karşılaştırmasıdır. Ortalama iki saat değişirken medyanın değişmemesi,
medyanın bu uç değere karşı daha sağlam olduğunu gösterir.